\subsection{The strict mixed cup product and the lower triangular interface}
\label{sec:cubic-strict-mixed-cup}

This section replaces the equal-character square in the source's lower
triangular construction by its exact mixed obstruction.  It proves the
cochain and local matrix interfaces.  It does not identify the resulting
changed-residual deformation space with the source's original
\path{Space O F 3}, whose \path{IsBar} clause still prescribes powers of
one character.

\paragraph{Conventions.}
Let \(k\) be a field of odd characteristic \(p\).  Let \(G\) be a profinite
group with a continuous involutive automorphism \(c\).  Let
\(\iota_v:G_v\to G\), for \(v\) in a finite set \(P\), be continuous local
maps.  We assume that \(P\) has a fixed-point-free involution \(v\mapsto
v^c\), with the same local group in a paired index and
\(\iota_{v^c}=c\iota_v\).  This is the source's literal doubled
\path{Frame.loc2} convention; no second copy is added after descent to
the real field.

Let \(\alpha:G\to k^\times\) be a continuous character with finite image,
put \(\beta=\alpha^c\), and put \(\omega=\alpha\beta\).
Here superscript \(c\) on a global cochain means precomposition by \(c\);
on a local family it permutes the paired indices.  In particular
\(\omega^c=\omega\).
All cochains below are continuous for the discrete coefficient module.
Our inhomogeneous differential is
\[
 d_\theta t(g)=(\theta(g)-1)t,\qquad
 d_\theta f(g,h)=f(g)+\theta(g)f(h)-f(gh).
 \tag{G90.1}
\]
Write
\[
 C_{\mathrm{str}}^n(G,k(\theta))
 =C^n(G,k(\theta))\oplus\bigoplus_{v\in P}C^{n-1}(G_v,k(\theta)),
 \qquad
 d(f,b)=(df,(\iota_v^*f-db_v)_v).
 \tag{G90.2}
\]
The second summand is zero in degree zero.

\begin{proposition}[AC-G90-cup: the strict quadratic obstruction]
\label{AC-G90-cup}
There is a natural homogeneous quadratic map
\[
 \mathcal Q:H^1_{\mathrm{str}}(G,k(\alpha))
       \longrightarrow H^2_{\mathrm{str}}(G,k(\omega))^{c=-1}.
 \tag{G90.3}
\]
It commutes with compatible restriction of the global and all local
groups and with finite coefficient-field extension.
If every \(H^0(G_v,k(\alpha))\) is zero, its domain is canonically the
strict kernel
\[
 \ker\!\left(H^1(G,k(\alpha))\longrightarrow
                 \bigoplus_v H^1(G_v,k(\alpha))\right).
 \tag{G90.4}
\]
\end{proposition}

\begin{proof}
Represent a degree-one relative cocycle by \((x,(a_v)_v)\), so that
\[
 d_\alpha x=0,\qquad
 \iota_v^*x=(\alpha-1)a_v.
\]
Set \(y=-x^c\) and \(b_v=-a_{v^c}\).  Then \(d_\beta y=0\) and
\(\iota_v^*y=(\beta-1)b_v\).
Use the ordered cup product
\[
 (y\smile x)(g,h)=y(g)\alpha(g)x(h),
 \qquad z_v^0(g)=\alpha(g)(\beta(g)-1)a_vb_v.
 \tag{G90.5}
\]
The cup product is a two-cocycle.  Direct substitution in (G90.1) gives
\(d_\omega z_v^0=-\iota_v^*(y\smile x)\).  Consequently
\[
 \mathcal Q(x,a)=
       [\,(-y\smile x,(z_v^0)_v)\,]
       \in H^2_{\mathrm{str}}(G,k(\omega)).
 \tag{G90.6}
\]

Here is a complete check that it descends to cohomology.  Replacing a
relative representative by a boundary means, for one \(t\in k\),
\[
 x'=x+d_\alpha t,\quad a_v'=a_v+t,\qquad
 y'=y-d_\beta t,\quad b_v'=b_v-t.
\]
Define the one-cochain
\[
 L_t(g)=t x(g)+\alpha(g)t y(g)
               +(\alpha(g)-\omega(g))t^2.
 \tag{G90.7}
\]
Expanding the cocycle formulas for \(x(gh)\) and \(y(gh)\) yields
\[
 d_\omega L_t=-y'\smile x'+y\smile x,\qquad
 (z_v^0)'-z_v^0=\iota_v^*L_t-d_\omega(t a_v).
 \tag{G90.8}
\]
Thus the two representatives of (G90.6) differ by the relative
boundary \(d(L_t,(t a_v)_v)\).

For the sign, \(y^c=-x\), and
\[
 d_\omega(xy)=-(x\smile y+y\smile x),\qquad
 \iota_v^*(xy)-z_v^0-(z_{v^c}^0)
                      =(1-\omega)a_vb_v.
 \tag{G90.9}
\]
The local identity follows by expanding
\((\alpha-1)(\beta-1)-\alpha(\beta-1)-\beta(\alpha-1)=1-\alpha\beta\);
also \(a_{v^c}b_{v^c}=a_vb_v\).
Hence the sum of the relative two-cocycle in (G90.6) and its conjugate is
exactly \(d(xy,(-a_vb_v)_v)\).  This proves the anti-invariance without
choosing a cup-product homotopy.

Every formula is quadratic in the relative representative.  Its
polarization is bilinear, and (G90.8) shows independence of
representatives, proving the asserted quadratic map on the quotient.
Restriction preserves the displayed formulas.  For a finite field
extension, choosing a basis identifies continuous cochains after
extension with the finite direct sum of the original cochain complex;
the same holds for local complexes and the \(c\)-eigenspaces.
Finally the long exact sequence of (G90.2) identifies its degree-one
cohomology with (G90.4) when the local degree-zero groups vanish.
\end{proof}

\begin{proposition}[AC-G90-fill: strict polarized lower triangular fillings]
\label{AC-G90-fill}
Put \(a=\alpha^{-1}\), \(d=\beta\), and
\[
 D(g)=\operatorname{diag}(a(g),1,d(g)),\qquad
 U(x,y,z)=
 \begin{pmatrix}1&0&0\\x&1&0\\z&y&1\end{pmatrix},\qquad
 J=\begin{pmatrix}0&0&1\\0&1&0\\1&0&0\end{pmatrix}.
 \tag{G90.10}
\]
For a representative \((x,(a_v))\) as above, \(\mathcal Q(x,a)=0\)
if and only if there is a continuous lower triangular representation
\[
 r(g)=U(x(g),-x^c(g),z(g))D(g)
 \tag{G90.11}
\]
which is conjugate to \(D|_{G_v}\) by
\(U(-a_v,-b_v,-C_v)\) for some constants \(C_v\)
at every local map and satisfies
\[
 r(cg)=Jr(g)^{-t}J.
 \tag{G90.12}
\]
When the local \(\alpha\)- and \(\beta\)-invariants vanish,
the adjacent entries of a strict local conjugating matrix are forced,
so this is equivalent to the unqualified strict local splitting
condition.
\end{proposition}

\begin{proof}
Matrix multiplication in (G90.11) gives precisely
\[
 d_\alpha x=0,\qquad d_\beta y=0,\qquad
 d_\omega z=-y\smile x.
 \tag{G90.13}
\]
For arbitrary constants \(A,B,C\), the unipotent coordinates of
\(U(A,B,C)D(g)U(A,B,C)^{-1}\) are
\[
 (1-\alpha)A,\quad (1-\beta)B,\quad
           (1-\omega)C+\alpha(\beta-1)AB.
 \tag{G90.14}
\]
Thus a vanishing relative obstruction is exactly the existence of
\(z\) and constants \(C_v\) such that
\[
 d_\omega z=-y\smile x,\qquad
 \iota_v^*z=z_v^0+d_\omega C_v.
 \tag{G90.15}
\]
The local conjugating matrix is \(U(-a_v,-b_v,-C_v)\).

One must impose polarity on the filling, not merely on its class.
The transformation
\[
 (z,C_v)\longmapsto
        (xy-z^c,\,-a_vb_v-C_{v^c})
 \tag{G90.16}
\]
is an affine involution of the set of solutions of (G90.15):
this is exactly (G90.9).  Since \(2\) is invertible, average a solution
with its transform.  The averaged solution has
\[
 z+z^c=xy,\qquad C_v+C_{v^c}=-a_vb_v.
 \tag{G90.17}
\]
The identities
\[
 JU(x,y,z)^{-t}J=U(-y,-x,xy-z),\qquad
 JD(g)^{-t}J=D(cg)
\]
prove (G90.12).  They also show that the two local conjugating
matrices at paired places are transforms of one another.
Conversely (G90.13) and (G90.14) give the relative boundary (G90.15)
from any representation and local constants in the statement.
\end{proof}

\begin{proposition}[AC-G90-Low: exact ordinary identities, normalization, and paths]
\label{AC-G90-Low}
The preceding strict local condition is equivalent to the source's
ordered ordinary matrix identity whenever, at each local map, there
is \(s_v\) for which \(a(s_v),1,d(s_v)\) are pairwise distinct:
\[
 (r(g_0)-a(g_0)I)(r(g_1)-I)(r(g_2)-d(g_2)I)=0
 \quad(g_0,g_1,g_2\in G_v).
 \tag{G90.18}
\]
This assertion also holds over any nonzero commutative \(k\)-algebra,
because the specified differences remain units.

Suppose moreover that there is \(g_\ast\in G\) with
\(d(g_\ast)=a(g_\ast)^{-1}\) and \(a(g_\ast)^2\ne1\).
Every polarized filling may be conjugated by a constant lower
unipotent matrix fixed by \(M\mapsto JM^{-t}J\), so that
\[
 r(g_\ast)-r(cg_\ast^{-1})\text{ is diagonal}.
 \tag{G90.19}
\]
A nonzero strict adjacent cohomology class remains nonzero after this
normalization.  If \(b\) is any strict \(\omega\)-cocycle with
\(b^c=-b\), then every normalized filling has the polynomial family
\[
 x_T=T x,\qquad y_T=T y,\qquad
 z_T=T^2z+(1-T)b,
 \tag{G90.20}
\]
which satisfies all the same matrix, local, polarization and
normalization equations.
\end{proposition}

\begin{proof}
Strict splitting immediately implies (G90.18), since diagonal
matrices commute and for each diagonal coordinate one of the three
factors is zero.  For the converse, write \(s=s_v\) and
\((d_1,d_2,d_3)=(a(s),1,d(s))\).  The spectral projectors
\[
 E_i=\prod_{h\ne i}\frac{r(s)-d_h I}{d_i-d_h}
\]
are lower triangular.  The matrix with columns \(E_i e_i\) is lower
unipotent and diagonalizes \(r(s)\).  After this constant conjugation,
specialize (G90.18) to \((s,g,s)\).  Its entries \(21,31,32\) are,
respectively, the corresponding entries of \(r(g)\) multiplied by
\[
 (1-d_1)(d_1-d_3),\quad
 (d_3-d_1)(d_1-d_3),\quad
 (d_3-d_1)(1-d_3).
\]
All are units.  Every \(r(g)\) is therefore diagonal in this same
basis.  This proves constant strict splitting, including over the
stated coefficient algebra.

For normalization put \(t=a(g_\ast)\), \(x_\ast=x(g_\ast)\), and
\(y_\ast=y(g_\ast)\).  Conjugate by
\[
 V=U(A,-A,-A^2/2),\qquad
 A=\frac{y_\ast-tx_\ast}{t-t^{-1}}.
 \tag{G90.21}
\]
The displayed \(V\) is fixed by \(M\mapsto JM^{-t}J\).
The new adjacent unipotent coordinates are
\(x+(1-\alpha)A\) and \(y-(1-\beta)A\).
Consequently their raw matrix entries at \(g_\ast\) differ by
\[
 t x_\ast-y_\ast+(t-t^{-1})A=0.
\]
For a lower matrix the star-transpose identity interchanges exactly
these two off-diagonal entries and fixes the long entry.  Thus
(G90.19) follows.  This conjugation changes \(x\) by a coboundary,
preserves strict local splitting and preserves polarity, so it cannot
kill a nonzero adjacent cohomology class.

For (G90.20), the mixed cup scales by \(T^2\), while \(d_\omega b=0\).
Both polarity and (G90.19) have the same homogeneous scaling.
If \(\iota_v^*b=d_\omega B_v\), the local primitives for the family are
\[
 T a_v,\quad T b_v,\quad T^2 C_v+(1-T)B_v.
\]
These verify (G90.15) for the family.  The paired constants can be
chosen anti-invariant by averaging; equivalently, apply (G90.16)
over \(k[T]\).  All coefficient functions remain locally constant.
\end{proof}

\paragraph{The actual cubic diagonal.}
For the residual diagonal
\[
 (\chi^j\gamma,\ 1,\ \chi^j\gamma^{-1}),\qquad
 \chi^c=\chi^{-1},\quad\gamma^c=\gamma,
\]
the coefficients are exactly
\[
 \alpha=\chi^{-j}\gamma^{-1},\qquad
 \beta=\chi^j\gamma^{-1},\qquad
 \omega=\gamma^{-2}.
 \tag{G90.22}
\]
At coefficient prime \(269\), an inertia witness for
\(\bar\gamma^2\ne1\) may be cubed to kill the order-three character
\(\bar\chi\).  Cubing cannot kill the inertia value of \(\bar\gamma^2\), whose
order divides \(268\).  This gives both the local separation in (G90.18) and an
eligible normalization element in (G90.21).  At a coefficient prime
where \(\chi\) is unramified, the original inertia witness already has
\(\chi=1\).  These statements retain the displayed residual
separation hypothesis; they do not cover coincident-character cases.

\begin{proposition}[AC-G90-quotient: the exact away-kernel comparison]
\label{AC-G90-quotient}
Let \(F\) be a number field, \(p\) odd, and let \(L/F\) be finite
Galois of degree prime to \(p\).  In one algebraic closure define
\(N_L\) to be the closed normal subgroup of \(G_F\) generated by
\[
 I_v\cap G_L\qquad(v\nmid p),
 \tag{G90.23}
\]
where all absolute valuations, hence all conjugates, are included.
Put \(\Gamma_L=G_F/N_L\).  Then there is an exact sequence
\[
 1\longrightarrow G_{L,S_p}\longrightarrow\Gamma_L
       \longrightarrow\operatorname{Gal}(L/F)\longrightarrow1.
 \tag{G90.24}
\]
Here \(G_{L,S_p}\) is the Galois group of the maximal extension of
\(L\) unramified outside the \(p\)-places and the infinite places.
At odd \(p\) the usual modified infinite local complexes contribute
zero to the compact-support comparison.

For a finite \(p\)-coefficient module \(M\) split by \(L\), the relative
global/p-local complex of \(\Gamma_L\) is the exact
\(\operatorname{Gal}(L/F)\)-invariant descent of the corresponding
complex over \(L\).  In particular it is legitimate to apply global
duality over \(L\) and then take these invariants.  This assertion does
not require inflation of degree-two cohomology into an unrestricted
absolute Galois group to be injective.

If \(L_0\subseteq L\) is Galois over \(F\) and already splits \(M\),
then inflation gives an isomorphism
\[
 H^1(\Gamma_{L_0},M)\ \simeq\ H^1(\Gamma_L,M).
 \tag{G90.25}
\]
Finally, for a cyclotomic \(\mathbb Z_p\)-tower \(F_n/F\), the subgroups
\(N_{L F_n}\), viewed in \(G_F\), equal \(N_L\).  The quotients for the
finite layers are the actual open subgroups of index \([F_n:F]\)
in (G90.24), with the usual compatible restriction and corestriction.
\end{proposition}

\begin{proof}
The family in (G90.23) is stable under conjugation by \(G_F\).
Its closed generated subgroup is therefore the same whether the
normal closure is taken inside \(G_L\) or inside \(G_F\).
Inside \(G_L\) these are exactly the inertia groups at every finite
place not above \(p\).  Killing their closed normal closure defines
the maximal extension unramified away from \(p\), proving (G90.24).

The Hochschild--Serre comparison for (G90.24) has no positive-degree
cohomology for its finite prime-to-\(p\) quotient: averaging is exact.
For the local term, group the places of \(L\) above a place \(v\) of
\(F\).  This is an induced family over the cosets of a decomposition
subgroup.  Taking invariants gives precisely the single complex
\(R\Gamma(F_v,M)\) by local Hochschild--Serre.  This proves the
relative-complex assertion, with every \(p\)-place counted once.
One can use functorial cochain resolutions for this argument; no
choice of a splitting of (G90.24) is being asserted.
The pairing and trace in compact-support global duality are natural
under these field automorphisms, so exact invariant descent preserves
the pairing.

For (G90.25), view both degree-one groups by inflation inside
\(H^1(G_F,M)\).  Their cocycles are precisely those vanishing on the
respective \(N\), since both \(N\)'s act trivially on \(M\).
On \(I_v\cap G_{L_0}\) such a cocycle is an additive homomorphism.
Its subgroup \(I_v\cap G_L\) has index prime to \(p\); a homomorphism
to the additive \(p\)-group \(M\) that vanishes there vanishes on the
whole group.  Thus a cocycle killing \(N_L\) also kills \(N_{L_0}\).
The reverse implication is immediate.  Coboundaries agree, proving
the isomorphism.

A cyclotomic \(\mathbb Z_p\)-extension is unramified at every finite
place away from \(p\).  Its absolute inertia groups at these places
are consequently unchanged, and
\[
 I_v\cap G_{L F_n}=I_v\cap G_L .
 \tag{G90.26}
\]
The whole family of such subgroups remains stable under \(G_F\).
The same closed-generation argument therefore identifies every
normal closure with \(N_L\).  It is contained in every \(G_{F_n}\);
the subgroup index and the restriction/corestriction assertions
follow.
\end{proof}

\paragraph{Native quotient and descent sign.}
The literal \path{HostFrameEx.GalLink.p_ker} kills the closed normal
closure of away-coefficient-prime inertia in
\(\ker\eta\cap\ker\bar\Psi\).  For the cubic residual tuple its
corresponding full kernel is
\(\ker\delta\cap\ker\bar\chi\cap\ker\bar\gamma\).
Choose \(L/F\), with \(F=H^+\), splitting that finite residual datum;
one may also adjoin \(\mu_p\), which does not change (G90.23), since
it is unramified away from \(p\).
The degree is prime to \(p\).
Because \(H/F\) is unramified at finite places, finite inertia already
lies in \(\ker\delta\).  Thus (G90.23) is exactly the proposed cubic
away-kernel, not an unrelated relaxed Selmer condition.

The anti eigenspace of the doubled complex over \(H\), with coefficient
\(\omega\), descends to the single-copy complex over \(F\) with
coefficient \(\omega\delta\).  Its compact-support dual character is
\[
 (\omega\delta)^{-1}\bar\epsilon
       =\delta\,\bar\psi^{\,2}\bar\epsilon^{\,2e+1}.
 \tag{G90.27}
\]
The character is even: both \(\delta\) and
\(\bar\epsilon^{2e+1}\) are odd and \(\bar\psi^2\) is even.
The cubic character has disappeared.  Equation (G90.25) explains
why a bound for its full degree-one cohomology may use its smaller
rational-character splitting field, although the finite cover used
to justify duality also splits \(\chi\).
The boundedness and eventual restriction consequences are supplied,
when their arithmetic hypotheses are met, in AC-G91; they are not
asserted by this algebraic comparison alone.

\paragraph{Exact source use and remaining interface.}
The pinned source \path{StrSMul.lean} defines \path{Frame.Low},
\path{Low.comb}, \path{Space.FPt}, and \path{Space.At}.
\path{LowerTriangularConjAntidiag.lean} constructs the lower
representation; its \path{ord3} identity already allows arbitrary
two endpoint entries.  Conversely
\path{OneAddMaximalIdealSubsetUnits.lean} obtains local splitting
from the ordered ordinary identity and one locally separated element.
The proof of (G90.18) supplies that converse directly for the cubic
endpoints.

The entire \path{ExistsNonzeroLowCocycle.lean} proof was read.
Its private \path{exists_nonzero_low_cocycle} builds a central entry
using a square of one adjacent cocycle.  Equations (G90.3)--(G90.17)
give its exact mixed-character replacement, conditional on the one
strict obstruction vanishing.  The subsequent
\path{preconnected_rank_three} argument uses only one nonzero
central strict cocycle and one nonzero adjacent Low point:
\path{Low.comb} connects points by weighted polynomial scaling.
Equation (G90.20) supplies precisely that scaling for the new
characters.  It does not require irreducibility of the whole
quadratic zero locus.

The source's \path{at_ex}, \path{at_uniq}, and \path{at_res} still
refer to its original universal space.  Its
\path{LeAdicClosure.NRep.IsBar} fixes the diagonal and local
characters to \((\Psi,1,\Psi^{-1})\).
The equations proved here verify the corresponding lower-matrix,
ordinary, polarity and normalization clauses for the cubic tuple;
they do not silently replace that \path{IsBar} clause or prove a
new universal \path{Space/At} theorem.  Nor do they establish a
nonzero actual strict adjacent class, the old frame inequalities,
or containment of a prescribed characteristic-\(p\) curve in
the automorphic ideal.

\paragraph{Source and geometric scope.}
All Lean files just cited are from commit
\path{7cbf7a0c32fca1badf05931879a37c3f8a2b03b4}; complete local
source files, hashes and acquisition receipts accompany this module.
No Lean build or proof of a changed declaration is claimed.
For the compact-support interpretation in AC-G90-quotient, the
published input is Milne, \emph{Arithmetic Duality Theorems},
II, Proposition~2.3 (the cone), Proposition~2.9 (the Galois
comparison), and Theorem~3.1(b) with Corollary~3.3 (the duality).
The new quotient and exact invariant-descent argument is given
above; a reproof of Milne's upstream duality theorem is not claimed.
In the already-audited SGA~1 dictionary, the finite residual
characters are finite etale local systems after removing their
ramified primes, and the lower matrices are their successive
extensions.  The mixed cup class is the obstruction to the
three-step extension with specified local splittings.  This paragraph
uses that existing dictionary and asserts no new archival reading.
